\begin{abstract}
\textbf{Background:} Spaceflight exposes biological systems to microgravity, inducing pronounced biomechanical stress and morphological reprogramming in plant cell walls. In NASA study OSD-615 (Advanced Plant EXperiment 03-1 / APEX-03-1), high-throughput glycome profiling across 155 monoclonal antibodies revealed extensive structural remodeling of non-cellulosic polysaccharides in \textit{Arabidopsis thaliana} seedling roots grown aboard the International Space Station (ISS) Veggie facility. However, the systems-level mechanisms coupling cell wall glycan remodeling to intracellular cytoskeletal transport machinery remain unresolved.

\textbf{Methods:} Here, we present an integrative systems biology framework that links continuous glycomic epitope matrices from OSD-615 to companion spaceflight transcriptomics (APEX-03-2 / OSD-218 and OSD-217) and interactome topologies. We employed sparse Partial Least Squares (sPLS) regression, Weighted Gene Co-expression Network Analysis (WGCNA), and protein--protein interaction (PPI) network modeling (STRING/AraNet) to correlate 155 glycan epitope vectors with motor protein complexes (kinesins, myosins), microtubule-associated proteins (MAP65, SPR1, CLASP), actin-regulatory machinery (ARP2/3, formins, profilins), and cellulose synthase complexes (CSCs). Furthermore, we developed a stochastic Monte Carlo vesicle transport model simulating secretory flux from Golgi stacks to the cell cortex under 1$g$ and microgravity conditions, and reviewed mass spectrometry workflows (EThcD, metabolic click-chemistry, lectin affinity) for mapping intracellular O-GlcNAcylation on motor complexes.

\textbf{Results:} Microgravity triggered significant alterations in cell wall glycomes, characterized by enhanced extractability of $\beta$-(1,4)-xylan epitopes (detected by CCRC-M138, CCRC-M139, CCRC-M140; $+1.8$ to $+2.4$ $\log_2\text{FC}$) and dynamic redistribution of arabinogalactan proteins (JIM13, JIM14) and xyloglucans (CCRC-M1, CCRC-M88). Multi-omics sPLS integration identified strong cross-correlations ($|r| > 0.75$) between xylan epitope abundance and transcriptional activation of secondary wall cellulose synthases (\textit{CESA4}, \textit{CESA7}), glycosyltransferases (\textit{IRX9}, \textit{IRX10}), and motor components (\textit{KIN12A}, \textit{MYA1}). Stochastic simulations demonstrated that microgravity-induced cortical microtubule disorientation and actin cable fragmentation decrease vesicle arrival efficiency at the cell plate by 28\%, directly altering matrix polysaccharide deposition rates.

\textbf{Conclusion:} These findings establish that microgravity-induced cell wall remodeling is mechanistically coupled to cytoskeletal motor dynamics and vesicular trafficking. We provide a FAIR-compliant open-access research repository, including an interactive 10-tab web dashboard, reproducible analysis pipelines, and full data dictionaries to support space agricultural engineering for extended lunar and Martian missions.
\end{abstract}
